\documentclass[11pt,a4paper]{article}
\usepackage[utf8]{inputenc}
\usepackage[T1]{fontenc}
\usepackage{lmodern}
\usepackage{tgpagella,tgheros}
\usepackage[margin=22mm,headheight=22pt,footskip=12mm]{geometry}
\PassOptionsToPackage{obeyspaces,spaces}{url}
\usepackage{microtype,graphicx,booktabs,tabularx,enumitem,listings,fancyhdr,titlesec,caption,xurl}
\usepackage[hidelinks,bookmarksnumbered=true]{hyperref}
\hypersetup{pdftitle={Food Court - Manual de usuario},pdfauthor={Ronald Martinez y Elias Fuentes},pdfsubject={Operación y evaluación de la API de la plazoleta de comidas},pdflang={es-CO}}
\renewcommand{\contentsname}{Contenido}
\renewcommand{\listtablename}{Tablas}
\renewcommand{\listfigurename}{Ilustraciones}
\renewcommand{\figurename}{Figura}
\renewcommand{\tablename}{Tabla}
\setcounter{tocdepth}{2}
\setlength{\parindent}{0pt}
\setlength{\parskip}{5pt}
\setlength{\tabcolsep}{5pt}
\renewcommand{\arraystretch}{1.26}
\setlist{leftmargin=*,itemsep=2pt,topsep=3pt}
\titleformat{\section}{\sffamily\bfseries\LARGE}{\thesection}{0.55em}{}
\titleformat{\subsection}{\sffamily\bfseries\Large}{\thesubsection}{0.55em}{}
\titleformat{\subsubsection}{\sffamily\bfseries\normalsize}{}{0pt}{}
\titlespacing*{\section}{0pt}{0pt}{9pt}
\titlespacing*{\subsection}{0pt}{9pt}{4pt}
\titlespacing*{\subsubsection}{0pt}{8pt}{2pt}
\captionsetup{font=small,labelfont=bf,justification=raggedright,singlelinecheck=false,hypcap=false}
\pagestyle{fancy}
\fancyhf{}
\fancyhead[L]{\small\sffamily\bfseries FOOD COURT}
\fancyhead[R]{\small\sffamily Manual de usuario $\cdot$ 1.0}
\fancyfoot[R]{\small\sffamily\bfseries\thepage}
\lstset{
  basicstyle=\ttfamily\small,
  frame=l,framerule=0.5pt,framesep=6pt,
  xleftmargin=8pt,xrightmargin=2pt,aboveskip=6pt,belowskip=7pt,
  breaklines=true,breakatwhitespace=false,columns=fullflexible,keepspaces=true,
  showstringspaces=false,upquote=true
}
\newcommand{\code}[1]{\nolinkurl{#1}}
\newcommand{\swagger}[1]{\href{http://localhost:#1/swagger-ui/index.html}{Swagger en #1}}
\newcommand{\route}[2]{\textbf{#1}\enspace\code{#2}}
\newcommand{\call}[3]{\route{#1}{#2}\enspace(puerto #3)}
\newcommand{\illustration}[3]{\begin{minipage}{\linewidth}\centering\includegraphics[width=#1,height=0.38\textheight,keepaspectratio]{assets/#2}\captionof{figure}{#3}\end{minipage}}
\newcolumntype{Y}{>{\raggedright\arraybackslash}X}
\newcolumntype{P}[1]{>{\raggedright\arraybackslash}p{#1}}
\raggedbottom
\begin{document}
\begin{titlepage}
\thispagestyle{empty}
\vspace*{19mm}
{\sffamily\bfseries\fontsize{42}{45}\selectfont Food Court\par}
\vspace{3mm}
{\sffamily\fontsize{27}{30}\selectfont Manual de usuario\par}
\vspace{9mm}
\rule{42mm}{1.8pt}\par
\vspace{7mm}
{\Large Operación de la API de la plazoleta de comidas\par}
\vfill
\begin{tabularx}{\linewidth}{@{}P{38mm}Y@{}}
\toprule
Integrantes & Ronald Martinez y Elias Fuentes \\
Versión documental & 1.0 \\
Fecha & 24 de noviembre de 2026 \\
Identificador & FC-MU-01 \\
\bottomrule
\end{tabularx}
\vspace{15mm}
\end{titlepage}
\clearpage
\setcounter{page}{2}
\begingroup
\small
\setlength{\parskip}{0pt}
\makeatletter
\renewcommand*{\l@section}[2]{\addvspace{3pt}\begingroup\bfseries\@dottedtocline{1}{0em}{2em}{#1}{#2}\endgroup}
\makeatother
\tableofcontents
\endgroup
\clearpage
\begingroup
\small
\setlength{\parskip}{0pt}
\listoftables
\listoffigures
\endgroup
\clearpage
\section{Introducción, objetivo y alcance}\label{sec:alcance}
En Food Court, el cliente elige platos de un restaurante y sigue su pedido hasta la entrega. El administrador registra los restaurantes y sus propietarios. Cada propietario mantiene el menú de sus locales, registra empleados y consulta los tiempos de entrega. El empleado atiende la cola de su restaurante y confirma la entrega con el PIN del cliente.

Este manual sirve para operar y evaluar las 18 historias de usuario de la entrega. Los procedimientos indican el rol, el puerto, los datos que debe ingresar y la respuesta que debe recibir. Los botones se indican en inglés, tal como aparecen en Swagger UI. El anexo \ref{sec:rutas} permite localizar una operación por historia de usuario.

Swagger UI es la interfaz de prueba incluida con las API. El alcance de la entrega es el backend, con cuatro microservicios y sus bases de datos. Los pagos y las aplicaciones web o móviles para clientes quedan fuera del alcance.

\subsection{Control documental}
\begin{tabularx}{\linewidth}{@{}P{18mm}P{34mm}Y@{}}
\toprule
Versión & Fecha & Descripción \\
\midrule
1.0 & 24/11/2026 & Manual inicial de la entrega final, con procedimientos por rol y un recorrido reproducible de pedido. \\
\bottomrule
\end{tabularx}

El manual toma como referencia el resumen público de ISO/IEC/IEEE 26514:2022 para organizar la información en tareas con resultados que el usuario pueda comprobar. Las fuentes se detallan en la sección \ref{sec:referencias}.

\clearpage
\section{Requisitos previos e instalación}\label{sec:instalacion}
Si el entorno ya está en marcha, abra Swagger en el navegador e inicie sesión con una cuenta del rol correspondiente. Para ejecutar el proyecto en su equipo, use Docker Compose. El glosario de la sección \ref{sec:glosario} reúne las herramientas y siglas del manual.

\subsection{Preparar una demostración con Docker}
\begin{enumerate}
\item Extraiga el ZIP y abra Bash en la carpeta que contiene \code{compose.yml}. Necesita Docker en ejecución, Compose 2.20.0 o posterior y conexión para descargar las imágenes y dependencias de la primera compilación. Use Bash, \code{curl} y \code{jq} para la semilla y el recorrido por terminal. En Windows puede usar Docker Desktop con una terminal Bash que tenga esas herramientas.
\item Compruebe que estén libres el puerto 5433 y los puertos del 8081 al 8084. Si todavía no existe \code{.env}, cópielo de la configuración de ejemplo.
\begin{lstlisting}
test -f .env || cp .env.example .env
\end{lstlisting}
\item Revise estos valores en \code{.env} antes de arrancar. \code{POSTGRES_USER} y \code{POSTGRES_PASSWORD} son las credenciales de PostgreSQL. \code{POSTGRES_DB} indica la base que Docker crea al inicializar el volumen. El proveedor local simula el aviso de pedido listo y \code{SEED_PIN} fija el PIN de la demostración.
\begin{lstlisting}
POSTGRES_USER=foodcourt
POSTGRES_PASSWORD=foodcourt
POSTGRES_DB=user_db

MESSAGING_PROVIDER=local
SEED_PIN=482913
\end{lstlisting}
\item Inicie el entorno. Docker ejecuta el SQL inicial al crear el volumen de PostgreSQL y cada servicio con base aplica sus migraciones con Flyway. Las rutas de los scripts están en la sección \ref{sec:postgres}.
\begin{lstlisting}
docker compose up -d --build
docker compose ps
\end{lstlisting}
\item Compruebe que los cuatro servicios publiquen OpenAPI. Si un servicio aún no responde, espere a que termine de arrancar y repita la consulta. Si el fallo persiste, revise la sección \ref{sec:soporte}.
\begin{lstlisting}
for port in 8081 8082 8083 8084; do
  curl -fsS "http://localhost:$port/v3/api-docs" -o /dev/null || break
  printf 'API disponible en %s\n' "$port"
done
\end{lstlisting}
\item En una instalación nueva, cargue la semilla una sola vez. Busque \code{Seed loaded.} al final y guarde los identificadores que muestra el script. En el primer arranque, Usuarios crea la cuenta de administrador con los datos de \code{.env}.
\begin{lstlisting}
bash scripts/seed.sh
\end{lstlisting}
\end{enumerate}

\clearpage
\subsection{Configurar PostgreSQL y ubicar los scripts}\label{sec:postgres}
Con Docker Compose, conéctese a PostgreSQL desde su equipo en \code{127.0.0.1:5433}. Los contenedores usan \code{postgres:5432}. Los servicios reciben las credenciales de \code{.env} y las URL de sus bases desde \code{compose.yml}.

\begin{minipage}{\linewidth}
\captionof{table}{Archivos de inicialización y datos de prueba}
\begin{tabularx}{\linewidth}{@{}P{74mm}Y@{}}
\toprule
Ruta desde la raíz del proyecto & Qué hace y cuándo se ejecuta \\
\midrule
\code{postgres/init-databases.sql} &
Crea \code{user_db}, \code{plaza_db} y \code{traceability_db} si faltan. Docker lo ejecuta al inicializar un volumen vacío. \\
\code{src/main/resources/db/migration/} en \code{user-service}, \code{plaza-service} y \code{traceability-service} &
Contiene los SQL que Flyway ejecuta cuando arranca el servicio correspondiente. Crean y actualizan las tablas, los índices y las restricciones de cada base. \\
\code{scripts/seed.sh} &
Carga los datos de demostración mediante la API. Ejecútelo una vez, desde la raíz, después de comprobar que los servicios respondan. \\
\bottomrule
\end{tabularx}
\end{minipage}

\textbf{Con PostgreSQL instalado en su equipo.} Use PostgreSQL 16 con soporte para ICU, \code{psql} y JDK 21. Desde la raíz del proyecto, ejecute el SQL como el usuario que será dueño de las bases. El ejemplo usa el rol \code{foodcourt} con permiso \code{CREATEDB} y el puerto habitual de una instalación nativa.
\begin{lstlisting}
psql -X -h 127.0.0.1 -p 5432 -U foodcourt -d postgres \
  -f postgres/init-databases.sql
\end{lstlisting}

En una terminal por servicio con base, cargue los valores de \code{.env} y exporte la conexión. El ejemplo muestra Usuarios. Para Plazoleta, use \code{plaza_db}. Para Trazabilidad, use \code{traceability_db}.
\begin{lstlisting}
set -a
source ./.env
set +a
export SPRING_DATASOURCE_URL=jdbc:postgresql://127.0.0.1:5432/user_db
export SPRING_DATASOURCE_USERNAME="$POSTGRES_USER"
export SPRING_DATASOURCE_PASSWORD="$POSTGRES_PASSWORD"
\end{lstlisting}

En la terminal de Mensajería basta con cargar \code{.env}. Desde la carpeta de cada servicio, ejecute este comando.
\begin{lstlisting}
./mvnw spring-boot:run
\end{lstlisting}

Cuando los cuatro servicios respondan, ejecute \code{bash scripts/seed.sh} desde la raíz. La guía independiente \code{IMPORTAR_BASE_DE_DATOS.md} detalla el arranque nativo.

\clearpage
\subsection{Servicios, acceso y cierre de la prueba}
\begin{minipage}{\linewidth}
\captionof{table}{Servicios y enlaces del entorno local}
\begin{tabularx}{\linewidth}{@{}P{33mm}P{34mm}P{31mm}Y@{}}
\toprule
Servicio & Entrada & Base & Función \\
\midrule
Usuarios & \swagger{8081} & \code{user_db} & Registro, sesión y cuentas. \\
Plazoleta & \swagger{8082} & \code{plaza_db} & Restaurantes, platos y pedidos. \\
Trazabilidad & \swagger{8083} & \code{traceability_db} & Historial y reporte de eficiencia. \\
Mensajería & \swagger{8084} & Sin base & Aviso de pedido listo. \\
\bottomrule
\end{tabularx}
\end{minipage}

Cada servicio publica su contrato en \code{/v3/api-docs}. Consulte \code{info.version} para conocer la versión del contrato. La plantilla del proyecto indica la revisión del código de cada servicio.

Para detener el entorno, ejecute \code{docker compose down} desde la raíz. Compose conserva el volumen. Si requiere una demostración nueva, detenga la instancia anterior y arranque con otro nombre de proyecto, por ejemplo \code{docker compose -p foodcourt-prueba up -d --build}. Use un nombre que todavía no tenga un volumen asociado y vuelva a ejecutar la semilla.

\section{Funciones y permisos}\label{sec:roles}
\begin{minipage}{\linewidth}
\captionof{table}{Acciones según el rol}
\begin{tabularx}{\linewidth}{@{}P{45mm}Y@{}}
\toprule
Rol en la API & Acciones \\
\midrule
Administrador \code{ADMINISTRATOR} & Registra propietarios y restaurantes. Puede consultar cuentas y la relación de propiedad de cualquier restaurante. \\
Propietario \code{OWNER} & Registra empleados en sus restaurantes, gestiona sus platos y consulta sus reportes de eficiencia. \\
Empleado \code{EMPLOYEE} & Lista pedidos de su restaurante, toma un pedido pendiente y completa sus pedidos con el PIN. \\
Cliente \code{CUSTOMER} & Registra su propia cuenta, consulta restaurantes y platos activos, pide, cancela si está pendiente y consulta su historial. \\
\bottomrule
\end{tabularx}
\end{minipage}

\subsection{Estados del pedido}
Un pedido comienza en \code{PENDING} (pendiente). Después pasa a \code{IN_PREPARATION} (en preparación), \code{READY} (listo) y \code{DELIVERED} (entregado). El cliente puede pasar de \code{PENDING} a \code{CANCELLED} (cancelado). Entregado y cancelado son estados finales. Mientras un pedido esté pendiente, en preparación o listo, ese cliente no puede abrir otro.

\clearpage
\section{Iniciar sesión y operar desde Swagger}\label{sec:sesion}
\subsection{Obtener un token}
\begin{enumerate}
\item Abra \swagger{8081}. En el grupo de autenticación, despliegue \call{POST}{/login}{8081} y pulse \textbf{Try it out}.
\item Sustituya el cuerpo por el correo y la contraseña de la cuenta. El ejemplo usa la cuenta de Mateo, incluida en la semilla. Pulse \textbf{Execute}.
\begin{lstlisting}
{"email":"mateo.vargas@example.com","password":"Password123"}
\end{lstlisting}
\item En \textbf{Server response}, confirme \code{200}. Copie el valor de \code{accessToken} sin las comillas. \code{userId} identifica la cuenta, \code{role} identifica el rol y \code{expiresIn} indica los segundos de vigencia. El entorno de ejemplo configura ocho horas.
\end{enumerate}

\illustration{0.94\linewidth}{swagger-login.png}{Operación de inicio de sesión en Swagger UI de Usuarios.}

Para consultar la cuenta que inició sesión, use \call{GET}{/users/{id}}{8081} con su \code{userId}. La respuesta incluye \code{id}, \code{role} y \code{restaurantId} cuando corresponde. El administrador puede consultar otras cuentas. Si pierde la contraseña, consulte al equipo responsable del entorno. La API no ofrece recuperación de contraseñas.

\clearpage
\subsection{Autorizar una llamada y cambiar de rol}
\begin{enumerate}
\item Abra Swagger en el puerto del servicio que ofrece la operación. Pulse \textbf{Authorize} e ingrese solo el valor de \code{accessToken} en \code{bearerAuth}. Swagger agrega el prefijo \code{Bearer}. Pulse \textbf{Authorize} dentro de la ventana y luego \textbf{Close}.
\item Abra la operación, pulse \textbf{Try it out}, ingrese los parámetros o el JSON y pulse \textbf{Execute}. Lea el código y el cuerpo en \textbf{Server response}. En Swagger, escriba el identificador real en el campo de ruta correspondiente. Los ejemplos usan los nombres \code{id}, \code{restaurantId} y \code{orderId}.
\item Para cambiar de rol, obtenga el token de la otra cuenta en \code{/login}. En cada servicio que vaya a usar, abra \textbf{Authorize}. Si ya hay un token cargado, pulse \textbf{Logout} y pegue el nuevo. Para dejar de usar una cuenta, retire su token de Swagger y de la terminal. Un JWT emitido conserva su vigencia hasta que expira.
\end{enumerate}

\illustration{0.78\linewidth}{swagger-authorize.png}{Cuadro de autorización Bearer de Swagger UI de Plazoleta.}

Si usa otro cliente HTTP, envíe \code{Authorization: Bearer <token>} y, cuando haya un cuerpo JSON, \code{Content-Type: application/json}. En la semilla, las cuentas de prueba comparten \code{Password123} salvo el administrador, cuya contraseña sale de \code{BOOTSTRAP_ADMIN_PASSWORD}. Las credenciales de ejemplo sirven para una instalación local de pruebas.

\clearpage
\section{Administrador}\label{sec:admin}
\subsection{Registrar un propietario y su restaurante}
Inicie sesión con la cuenta que figura en \code{BOOTSTRAP_ADMIN_EMAIL}. En la configuración de ejemplo, el correo es \code{admin@example.com}. Autorice el mismo token en Usuarios y en Plazoleta.

\textbf{1. Crear la cuenta propietaria (HU 1).} Ejecute \call{POST}{/owners}{8081}. Use un correo y un documento que no estén registrados. El propietario debe tener al menos 18 años. Escriba el documento y el teléfono como cadenas de dígitos. El teléfono admite el signo \code{+} al inicio. Use una contraseña de 8 a 72 caracteres y una fecha \code{AAAA-MM-DD}.
\begin{lstlisting}
{
  "firstName":"Nora", "lastName":"Rios",
  "identityDocument":"899100201", "phone":"+573001002201",
  "birthDate":"1992-04-20", "email":"nora.rios@example.com",
  "password":"Password123"
}
\end{lstlisting}
Confirme \code{201} y guarde el \code{id}. El rol se asigna automáticamente. Si obtiene \code{409}, revise si se repitió el documento o el correo.

\textbf{2. Registrar el restaurante (HU 2).} Ejecute \call{POST}{/restaurants}{8082}. Sustituya \code{123} por el \code{id} de Nora que devolvió la operación anterior. Escriba el NIT como una cadena de solo dígitos, con hasta 15 caracteres. Use un nombre que contenga letras, un teléfono de hasta 13 caracteres y una URL de logo \code{http} o \code{https}.
\begin{lstlisting}
{
  "name":"Cocina Norte", "nit":"901288203",
  "address":"Calle 20 # 12-34", "phone":"+573001002202",
  "logoUrl":"https://example.com/cocina-norte.png", "ownerId":123
}
\end{lstlisting}
Confirme \code{201}. Guarde el \code{id} del restaurante y verifique que \code{ownerId} corresponde a Nora. El nombre y el NIT deben ser únicos. Un mismo propietario puede gestionar más de un restaurante. La URL de ejemplo ilustra el formato esperado.

\textbf{3. Verificar la asignación.} Ejecute \call{GET}{/restaurants/{id}}{8082} con el id recién creado. La respuesta \code{200} contiene \code{id} y \code{ownerId}. El administrador puede consultar cualquier restaurante. El propietario puede consultar los que tiene a cargo.

\illustration{0.89\linewidth}{swagger-restaurant.png}{Creación de Cocina Norte con el identificador del propietario.}

\clearpage
\section{Propietario}\label{sec:owner}
Para trabajar con el restaurante que creó, inicie sesión con Nora y use el identificador de Cocina Norte. Si va a usar los datos de prueba, inicie sesión con \code{carlos.ramirez@example.com}. El valor \code{restaurantId: 1} de los dos ejemplos corresponde a La Parrilla de Carlos. Si usa la cuenta de Nora, reemplácelo por el identificador de Cocina Norte.

\subsection{Registrar empleados (HU 6)}
Ejecute \call{POST}{/employees}{8081}. Use \code{roleId: 3}, el id de un restaurante propio y datos de una cuenta nueva. El empleado queda vinculado al restaurante y, después de iniciar sesión, puede gestionar sus pedidos.
\begin{lstlisting}
{
  "firstName":"Elena", "lastName":"Rios",
  "identityDocument":"899100202", "phone":"+573001002203",
  "email":"elena.rios@example.com", "password":"Password123",
  "roleId":3, "restaurantId":1
}
\end{lstlisting}
Confirme que la API responde \code{201}, asigna el rol \code{EMPLOYEE} y devuelve el restaurante en \code{restaurantId}. Si la API devuelve \code{403}, revise que su cuenta sea propietaria de ese local.

\subsection{Crear un plato (HU 3)}
Ejecute \call{POST}{/dishes}{8082}. El precio es un entero positivo. Incluya la categoría, una descripción y una URL de imagen con \code{http} o \code{https}. Use una URL real si desea que un cliente externo pueda cargar la imagen.
\begin{lstlisting}
{
  "name":"Arepa de maiz", "price":12000,
  "description":"Arepa asada con queso fresco",
  "imageUrl":"https://example.com/arepa.png",
  "category":"Entradas", "restaurantId":1
}
\end{lstlisting}
Confirme \code{201} y \code{active: true}. Conserve el \code{id} del plato para editarlo. El menú del cliente incluye el plato cuando está activo.

\subsection{Editar el precio y la descripción (HU 4)}
Ejecute \call{PUT}{/dishes/{id}}{8082} con el id del plato propio. El cuerpo admite los dos campos siguientes. Confirme \code{200} y los nuevos valores en el resultado.
\begin{lstlisting}
{"price":13000,"description":"Arepa asada con queso y hogao"}
\end{lstlisting}

\clearpage
\subsection{Desactivar o volver a ofrecer un plato (HU 7)}
Ejecute \call{PUT}{/dishes/{id}/active}{8082} con uno de estos cuerpos, según la acción. Use el id del plato creado en la sección anterior. Confirme una respuesta \code{200} y compruebe el valor de \code{active}.
\begin{lstlisting}
{"active":false}
\end{lstlisting}
\begin{lstlisting}
{"active":true}
\end{lstlisting}
Para comprobar el efecto, cambie al token de cliente y consulte \code{GET /restaurants/{restaurantId}/dishes}. Busque el plato por categoría o revise todas las páginas. El plato desactivado queda fuera del menú activo y de su recuento. Vuelve al menú al reactivarlo. El sistema conserva el registro del plato.

\illustration{0.89\linewidth}{swagger-dish-active.png}{La respuesta muestra el plato con el campo active en false.}

\clearpage
\subsection{Consultar eficiencia (HU 18)}\label{sec:eficiencia}
Inicie sesión como propietario y abra \swagger{8083}. Ejecute \call{GET}{/restaurants/{restaurantId}/efficiency}{8083} con su restaurante. Puede enviar \code{page=0} y \code{size=10}. La primera página es 0 y el tamaño permitido va de 1 a 50.

\begin{itemize}
\item \code{completedOrders} cuenta los pedidos con traza desde \code{IN_PREPARATION} hasta \code{DELIVERED}. \code{averageDurationSeconds} indica la duración media de esos pedidos.
\item \code{orders.items} muestra el id del pedido, el empleado, las dos marcas de tiempo y \code{durationSeconds}. Los más largos aparecen primero. El promedio se expresa en segundos enteros. \code{slowerThanAverage} señala una duración estrictamente superior al promedio calculado.
\item \code{employeeRanking} resume todos los pedidos terminados, aunque cambie la página. El puesto 1 corresponde al menor promedio. Si los promedios empatan, primero aparece el empleado con el identificador más bajo.
\end{itemize}

Inmediatamente después de la semilla, el restaurante 1 devuelve \code{completedOrders: 0}, \code{averageDurationSeconds: null}, una página sin elementos y un ranking vacío. Complete la entrega del pedido de Ronald para obtener el primer pedido completado. El reporte muestra el empleado asignado y los segundos transcurridos entre la asignación y la entrega. Consulte el recorrido de la sección \ref{sec:demo}.

\illustration{0.94\linewidth}{swagger-efficiency.png}{Reporte de La Parrilla de Carlos después de la entrega. Pedro es el empleado de id 5.}

\clearpage
\section{Cliente}\label{sec:cliente}
\subsection{Crear una cuenta (HU 8)}
En \swagger{8081}, ejecute \route{POST}{/customers}. El registro es público. Use \code{roleId: 4}, un correo sin registrar y un documento nuevo de hasta 20 dígitos. Ingrese una contraseña de 8 a 72 caracteres y un teléfono de 7 a 13 caracteres. Para SMS real, use un teléfono autorizado con prefijo internacional y signo \code{+}.
\begin{lstlisting}
{
  "firstName":"Irene", "lastName":"Rios",
  "identityDocument":"899100203", "phone":"+573001002204",
  "email":"irene.rios@example.com", "password":"Password123",
  "roleId":4
}
\end{lstlisting}
Confirme \code{201} y el rol \code{CUSTOMER}. Inicie sesión con el nuevo correo. Si usó la semilla, también puede pedir con la cuenta de Mateo, que queda sin un pedido activo.

\subsection{Buscar un restaurante y consultar su menú (HU 9 y 10)}
\begin{enumerate}
\item Autorice el token en \swagger{8082} y ejecute \route{GET}{/restaurants} con \code{page=0} y \code{size=10}. La lista se ordena por nombre e incluye id, nombre y URL del logo. Guarde el id del restaurante.
\item Ejecute \route{GET}{/restaurants/{restaurantId}/dishes}. Use el id 1 para La Parrilla de la semilla. En \code{category}, escriba \code{Carnes} para filtrar o deje el campo vacío para ver todos los platos activos. Se ordenan por nombre y, en caso de empate, por id.
\item Anote los identificadores de los platos en \code{items}. Consulte \code{totalItems} y \code{totalPages}. Aumente \code{page} para avanzar y use un \code{size} de 1 a 50. Si la página está vacía, revise el filtro y la posición de la página.
\end{enumerate}

\illustration{0.83\linewidth}{swagger-menu.png}{Menú de Carnes con un plato por página y dos resultados en total.}

\subsection{Realizar un pedido (HU 11)}\label{sec:pedir}
Ejecute \call{POST}{/orders}{8082} con el token de Irene o Mateo. Este ejemplo elige el Churrasco de La Parrilla (plato 1 en la semilla). Si cambió de instalación, use los identificadores que devuelve el menú.
\begin{lstlisting}
{"restaurantId":1,"lines":[{"dishId":1,"quantity":2}]}
\end{lstlisting}
El pedido admite de 1 a 20 líneas, cada una con un plato activo distinto del mismo restaurante y una cantidad entera entre 1 y 50. Confirme \code{201} y \code{status: PENDING}. Guarde el \code{id} del pedido. Si el cliente ya tiene un pedido activo, la operación devuelve \code{409}.

\clearpage
\subsection{Cancelar un pedido pendiente (HU 16)}
Con el token del cliente que creó el pedido, ejecute \call{POST}{/orders/{orderId}/cancellation}{8082}. Ingrese el id y ejecute con el cuerpo vacío. Si el pedido sigue pendiente, la API devuelve \code{200} con \code{status: CANCELLED}. El cliente ya puede hacer otro pedido. Si el pedido ya cambió de estado, recibe \code{409}. La API devuelve \code{404} si el pedido no existe o pertenece a otro cliente.

Para probar la cancelación, use el pedido nuevo de la sección \ref{sec:pedir} antes de que un empleado lo tome. Si cargó la semilla, el pedido 2 de Lucía ya está cancelado. Puede consultarlo con su token en trazabilidad para ver el resultado guardado.

\subsection{Consultar los cambios de estado (HU 17)}\label{sec:historial}
Abra \swagger{8083}, autorice el token del cliente que hizo el pedido y ejecute \call{GET}{/orders/{orderId}/trace}{8083}. Use el id del pedido, aunque ya esté entregado o cancelado.

La respuesta \code{200} contiene \code{orderId}, \code{customerId}, \code{restaurantId} y \code{changes}, en orden cronológico. Cada cambio muestra \code{previousStatus}, \code{newStatus}, \code{occurredAt} y \code{secondsInStatus}. El primer cambio tiene estado anterior nulo. En el último, \code{secondsInStatus} es nulo porque no hay un cambio posterior. Los tiempos se muestran como marcas ISO 8601. Compare los estados con la siguiente secuencia del pedido de Ronald al terminar la semilla.
\begin{lstlisting}
null           -> PENDING
PENDING        -> IN_PREPARATION
\end{lstlisting}
La API devuelve \code{404} si el historial del pedido no existe o corresponde a otro cliente. Por ejemplo, Mateo recibe \code{404} al intentar consultar el pedido de Ronald.

\illustration{0.89\linewidth}{swagger-trace.png}{Historial del pedido de Ronald después de la entrega.}

\subsection{Recibir el PIN y reclamar el pedido}
Cuando el empleado marque el pedido listo, el proveedor de SMS configurado recibe la solicitud de envío al número registrado por el cliente. Entregue el PIN al empleado al reclamar el pedido. En el entorno local de este manual, la mensajería se simula y el evaluador utiliza \code{482913}. Después de la entrega, el cliente puede ver \code{DELIVERED} en el historial.

\clearpage
\section{Empleado}\label{sec:empleado}
\subsection{Consultar la cola de pedidos (HU 12)}
Inicie sesión con su cuenta de empleado y cargue el token en \swagger{8082}. Para la demostración, use \code{pedro.martinez@example.com}. Ejecute \call{GET}{/orders}{8082} con \code{status=PENDING}, \code{page=0} y \code{size=10}. El parámetro \code{status} es obligatorio y utiliza uno de los cinco códigos de la sección \ref{sec:roles}.

La lista se limita al restaurante registrado en su cuenta. Cada resultado muestra el id, el cliente, el restaurante, el estado, el empleado asignado y las líneas. Un pedido pendiente tiene \code{assignedEmployeeId: null}. También puede consultar otros estados del mismo restaurante. Ver el pedido de otro empleado en la cola no concede permiso para completarlo.

\illustration{0.94\linewidth}{swagger-orders.png}{Cola de Pedro filtrada por pedidos en preparación. El pedido 1 ya tiene empleado asignado.}

\subsection{Tomar un pedido pendiente (HU 13)}
1. Elija un pedido de la cola \code{PENDING} y anote su identificador. Si quiere completar un pedido de Mateo o Irene, déjelo pendiente hasta que un empleado lo tome.

2. Ejecute \call{POST}{/orders/{orderId}/assignment}{8082} con ese id y cuerpo vacío. Verifique \code{200}, \code{status: IN_PREPARATION} y su id en \code{assignedEmployeeId}. Consulte de nuevo la cola para comprobar el cambio. Si otro empleado ya lo tomó, la llamada devuelve \code{409}. Si el pedido es de otro restaurante, devuelve \code{404}.

\clearpage
\subsection{Marcar un pedido como listo y solicitar el aviso (HU 14)}\label{sec:listo}
Ejecute \call{POST}{/orders/{orderId}/ready}{8082} con cuerpo vacío. Use el token del empleado asignado y un pedido en \code{IN_PREPARATION}. Para el recorrido de la semilla, Pedro ya es el empleado asignado al pedido 1 de Ronald.

Si la operación se completa, la API devuelve \code{200} con \code{status: READY}. La API devuelve el pedido sin el PIN ni su hash. Plazoleta consulta el teléfono del cliente y pide a Mensajería que envíe el PIN. Con el proveedor \code{local} se acepta el envío de prueba sin entregar un SMS. Si recibe un error del proveedor o de trazabilidad, revise la cola antes de reintentar. Plazoleta intenta restaurar el estado anterior si falla el aviso o el registro de la traza.

\illustration{0.86\linewidth}{swagger-ready.png}{Respuesta de Plazoleta después de marcar el pedido 1 como listo.}

\subsection{Confirmar la entrega (HU 15)}
Para un pedido \code{READY}, ejecute \call{POST}{/orders/{orderId}/delivery}{8082} como el mismo empleado. Use el PIN de seis dígitos indicado por el cliente. Escríbalo como texto para conservar ceros al inicio. El siguiente valor corresponde a la configuración de demostración.
\begin{lstlisting}
{"pin":"482913"}
\end{lstlisting}
Confirme \code{200} y \code{status: DELIVERED}. Si el PIN no coincide, la API devuelve \code{409} y el pedido sigue listo. Si el PIN tiene un formato inválido, devuelve \code{400}. El pedido ya entregado no admite una segunda entrega.

\illustration{0.94\linewidth}{swagger-delivery.png}{Confirmación de entrega del pedido 1 tras ingresar el PIN de demostración.}

\clearpage
\section{Recorrido reproducible de evaluación}\label{sec:demo}
Ejecute este recorrido inmediatamente después de cargar la semilla con los valores de la sección \ref{sec:instalacion}. En una primera carga sobre bases vacías, La Parrilla y el pedido de Ronald tienen el identificador 1 en sus respectivas listas. Si el identificador difiere, consulte la salida de la semilla y la cola de Pedro. Los pasos pueden ejecutarse en Swagger o por terminal según el anexo \ref{sec:terminal}.

\begin{minipage}{\linewidth}
\captionof{table}{Cuentas para el recorrido}
\begin{tabularx}{\linewidth}{@{}P{30mm}P{63mm}Y@{}}
\toprule
Persona & Correo & Uso \\
\midrule
Carlos & \code{carlos.ramirez@example.com} & Propietario de La Parrilla y Sushi Zen 88. \\
Pedro & \code{pedro.martinez@example.com} & Empleado de La Parrilla. \\
Ronald & \code{ronald.martinez@example.com} & Cliente del pedido que se completará. \\
Mateo & \code{mateo.vargas@example.com} & Cliente sin pedido activo al terminar la semilla. \\
Lucía & \code{lucia.navarro@example.com} & Cliente con un pedido cancelado y otro pendiente. \\
\bottomrule
\end{tabularx}
\end{minipage}

\begin{enumerate}
\item Como Carlos, consulte \call{GET}{/restaurants/1/efficiency}{8083}. Compruebe \code{completedOrders: 0}.
\item Como Pedro, consulte \call{GET}{/orders?status=IN_PREPARATION}{8082}. La cola contiene el pedido 1 asignado a Pedro. Con el mismo token, ejecute \call{POST}{/orders/1/ready}{8082}. Debe quedar \code{READY}.
\item Intente entregar con \code{"pin":"000000"}. Debe recibir \code{409}. Repita con \code{"pin":"482913"}. La respuesta correcta devuelve \code{DELIVERED}.
\item Como Ronald, consulte \call{GET}{/orders/1/trace}{8083}. La secuencia esperada es \code{PENDING}, \code{IN_PREPARATION}, \code{READY} y \code{DELIVERED}. La última duración debe ser nula. Como Mateo, la misma consulta debe responder \code{404}.
\item Vuelva a usar el token de Carlos y consulte \code{/restaurants/1/efficiency}. Debe mostrar un pedido completado, con Pedro en el puesto 1. Con un solo pedido, \code{slowerThanAverage} es \code{false}. Los segundos dependen del momento en que se cargó la semilla y se hizo la entrega.
\item Para probar la cancelación, cree un pedido con la cuenta de Mateo según la sección \ref{sec:pedir}. Cancélelo mientras siga pendiente y consulte su traza. Debe mostrar \code{PENDING} y \code{CANCELLED}.
\end{enumerate}

La guía de datos \code{docs/testing/seed-data.md} lista las cuentas y platos restantes. Para repetir el recorrido desde el mismo estado inicial, use una instancia de pruebas nueva y vuelva a cargar la semilla.

\clearpage
\section{Mensajería y restricciones de operación}\label{sec:sms}
Al marcar un pedido como listo con \code{/orders/{orderId}/ready}, Plazoleta solicita el aviso. El anexo \ref{sec:rutas} indica las rutas de Mensajería y Trazabilidad que puede consultar en Swagger para diagnosticar fallos.

\begin{minipage}{\linewidth}
\captionof{table}{Configuración del proveedor de SMS}
\begin{tabularx}{\linewidth}{@{}P{27mm}Y@{}}
\toprule
Proveedor & Configuración y resultado \\
\midrule
\code{local} & Para la evaluación. El proveedor acepta la solicitud de prueba sin enviar un SMS. Use \code{SEED_PIN=482913} en Compose para conocer el PIN de prueba. \\
\code{textbee} & Opción predeterminada. Solicita un SMS real. Configure \code{TEXTBEE_API_KEY} y mantenga un Android registrado en línea, con una SIM que pueda enviar SMS. Use \code{TEXTBEE_DEVICE_ID} si quiere elegir el dispositivo remitente. \\
\bottomrule
\end{tabularx}
\end{minipage}

\code{MESSAGING_PROVIDER} se aplica en Mensajería. Sus valores admitidos son \code{local} y \code{textbee}. Un valor desconocido impide el inicio de ese servicio. \code{SEED_PIN} se aplica en Plazoleta, donde Compose lo pasa como \code{APP_PIN_FIXED}. Un valor de seis dígitos fija el PIN de los pedidos que se marquen listos. Si queda vacío, Plazoleta genera un PIN aleatorio para cada pedido.

\code{scripts/seed.sh} funciona con \code{SEED_PIN} configurado o vacío. Deja el pedido de Ronald en \code{IN_PREPARATION}. Para completar la entrega con el proveedor local, inicie Plazoleta con un PIN fijo conocido. Una configuración de PIN fijo inválida impide que Plazoleta arranque. Si ejecuta Plazoleta directamente con Java, exporte \code{SEED_PIN}. Si también define \code{APP_PIN_FIXED}, este último tiene prioridad.

Para probar un SMS real, seleccione \code{textbee} y configure sus credenciales en \code{.env}. Use \code{TEXTBEE_BASE_URL} para cambiar la URL predeterminada si corresponde. Para generar PIN aleatorios, deje \code{SEED_PIN} vacío. Si cambia el valor con el entorno en marcha, aplique el cambio con \code{docker compose up -d plaza-service}. Los pedidos que ya están listos conservan su PIN. Cree un cliente con un número propio autorizado, haga un pedido y siga el recorrido de empleado. Confirme que el mensaje llegó y use el PIN recibido.

\clearpage
\subsection{Del aviso a la entrega}
\begin{enumerate}
\item El empleado asignado ejecuta \call{POST}{/orders/{orderId}/ready}{8082}. Plazoleta consulta el número en \call{GET}{/users/{id}/phone}{8081}, genera el PIN y guarda su hash al pasar de \code{IN_PREPARATION} a \code{READY}.
\item Plazoleta llama a \call{POST}{/notifications/order-ready}{8084} con el identificador del pedido, el restaurante, el teléfono y el PIN. En \code{local}, Mensajería acepta la solicitud sin llamar a un proveedor externo. En \code{textbee}, la envía a \code{POST /gateway/send-sms}. Si el proveedor acepta, Mensajería devuelve \code{202} con \code{channel: SMS}. Ese código confirma aceptación y el cliente comprueba la recepción en su teléfono.
\item Una vez aceptado el aviso, Plazoleta registra \code{READY} mediante \call{POST}{/order-status-changes}{8083}. El empleado recibe \code{200} de Plazoleta. El cliente consulta el estado en \call{GET}{/orders/{orderId}/trace}{8083}.
\item Para entregar, el mismo empleado ejecuta \call{POST}{/orders/{orderId}/delivery}{8082} con \code{pin} como texto de seis dígitos. Plazoleta compara el PIN con el hash. Si coincide, registra \code{DELIVERED} y responde \code{200}. La entrega no solicita otro aviso.
\end{enumerate}

\illustration{0.87\linewidth}{swagger-notification.png}{Respuesta de Mensajería en modo local. HTTP 202 indica que aceptó la solicitud.}

Las rutas de integración se incluyen para identificar dónde ocurrió un error. Para cambiar el estado de un pedido, use las rutas de Plazoleta. Los ejemplos de este manual usan \code{local}. Mantenga las claves, los tokens y los PIN fuera de reportes y capturas.

\clearpage
\section{Solución de problemas y soporte}\label{sec:soporte}
En Swagger, lea primero el código de la respuesta y su \code{detail}. Si hay errores de campos, compare el cuerpo enviado con el JSON del procedimiento. Para fallos entre servicios, lea el registro de los contenedores correspondientes.
\begin{lstlisting}
docker compose ps
docker compose logs --tail=80 user-service plaza-service
docker compose logs --tail=80 traceability-service messaging-service
\end{lstlisting}

\subsection{Diagnóstico por síntoma}
\begin{minipage}{\linewidth}
\captionof{table}{Respuestas frecuentes y acciones de recuperación}
\begin{tabularx}{\linewidth}{@{}P{43mm}Y@{}}
\toprule
Síntoma & Qué revisar \\
\midrule
Swagger no abre & Confirme el puerto, el estado de los contenedores y que \code{/v3/api-docs} responda. Un contenedor recién iniciado puede tardar en publicar la API. \\
\code{400} Solicitud inválida & Verifique JSON, campos requeridos, tipo y formato. En listados revise \code{page}, \code{size}, el estado en inglés y los id positivos. Para entrega, envíe el PIN como texto de seis dígitos. \\
\code{401} Autenticación & Revise credenciales al iniciar sesión. En rutas protegidas, obtenga un token nuevo y cárguelo en el servicio correcto. Compruebe que los servicios compartan \code{JWT_SECRET}. \\
\code{403} Permisos & Compruebe el rol de la cuenta y, donde aplique, la propiedad del restaurante. Cambie de token en Swagger si conservó el de una cuenta anterior. \\
\code{404} Recurso fuera de alcance & Confirme el id y que consulta el puerto indicado. En pedidos o trazas ajenas, el servicio también puede responder como si el recurso no existiera. \\
\code{409} Conflicto & Revise duplicados, pedido activo, estado del pedido, empleado asignado o PIN incorrecto. Consulte la cola o la traza antes de repetir. \\
\code{415} Formato de contenido & Envíe \code{Content-Type: application/json} cuando la operación requiere un cuerpo JSON. \\
\code{502} o \code{503} Integración & Revise el servicio dependiente y su configuración. Si se rechazó un aviso o una traza, compruebe el estado registrado del pedido antes de reintentar. \\
Semilla interrumpida & Lea el primer error de la salida. Una carga parcial puede haber creado cuentas o pedidos. Use una instancia de pruebas nueva para repetirla. \\
\bottomrule
\end{tabularx}
\end{minipage}

\clearpage
\subsection{Recuperar un pedido después de un error}
\begin{enumerate}
\item Si \code{/ready} respondió \code{502}, revise los registros de Mensajería. Si usa textbee, compruebe su clave. Si obtiene \code{503}, compruebe que los servicios dependientes estén disponibles. Plazoleta intenta restaurar \code{IN_PREPARATION} cuando falla el aviso o el registro de la traza.
\item Consulte la cola de Plazoleta con el token de empleado y el historial con el token de cliente. Si el pedido sigue en preparación, solucione el problema de integración y vuelva a marcarlo listo. Si ya aparece listo, confirme si se recibió el aviso antes de repetir acciones.
\item Si \code{/delivery} devolvió \code{409} por PIN incorrecto, el pedido permanece listo. Confirme los seis dígitos con el cliente o use el PIN fijo de la demostración. Conserve los ceros iniciales al introducirlo.
\end{enumerate}

\illustration{0.90\linewidth}{swagger-pin-error.png}{Respuesta de Plazoleta al entregar un pedido listo con un PIN que no coincide.}

\clearpage
\subsection{Preguntas frecuentes}
\textbf{¿Por qué el reporte está vacío?} El reporte requiere al menos un pedido con traza de preparación y entrega para ese restaurante. Los pendientes y cancelados quedan fuera del cálculo.

\textbf{¿Por qué no llegó el SMS?} En \code{local} se simula el envío. Con un proveedor externo, confirme el teléfono con prefijo, la configuración del proveedor y el estado real del pedido. Si Plazoleta respondió con error de integración, consulte los registros y la cola luego de restablecer el servicio.

\textbf{¿Cómo vuelvo a empezar?} Detenga el entorno y use otro nombre de proyecto para obtener un volumen nuevo, como se explica en la sección \ref{sec:instalacion}. Cada comando posterior de Compose debe usar el mismo nombre de proyecto.

\subsection{Reportar un incidente}
Envíe el reporte a Ronald Martinez y Elias Fuentes por el canal acordado. Incluya la hora, el rol, el puerto, la ruta y el identificador del pedido. Anote los pasos para reproducirlo, el resultado esperado, el código de estado HTTP y el campo \code{detail} de la respuesta. Adjunte únicamente los mensajes del registro que describen el fallo, sin credenciales, teléfonos ni PIN.

\textbf{Ejemplo de reporte.} Como empleado, abrí el pedido 1 en la cola \code{READY} y ejecuté \code{POST /orders/1/delivery} en el puerto 8082. Esperaba \code{200} y \code{DELIVERED}. Recibí \code{409} con \code{The security PIN does not match}. Volví a consultar la cola y el pedido seguía en \code{READY}.

\clearpage
\section{Términos y definiciones}\label{sec:glosario}
\begin{tabularx}{\linewidth}{@{}P{34mm}Y@{}}
\toprule
Término & Significado en Food Court \\
\midrule
API REST & Interfaz que recibe solicitudes HTTP. Una operación se identifica por su método y su ruta. \\
Bearer / JWT & El JWT es el token temporal de la sesión. Se envía con el esquema Bearer en el encabezado \code{Authorization}. \\
Docker Compose & Herramienta que arranca PostgreSQL y los microservicios con la configuración del proyecto. \\
Flyway & Herramienta que instala y registra las versiones del esquema en cada base. \\
ID & Identificador numérico de una cuenta, restaurante, plato o pedido. Se obtiene de la respuesta correspondiente. \\
JSON & Formato de los cuerpos y respuestas. Los textos van entre comillas dobles. \\
OpenAPI & Descripción de rutas, parámetros y formatos que publica cada servicio. Swagger UI la presenta en el navegador. \\
PIN & Código de seis dígitos que el cliente muestra para retirar un pedido listo. \\
Semilla & Datos que \code{scripts/seed.sh} crea mediante solicitudes a las API para repetir pruebas. \\
Trazabilidad & Historial cronológico de los estados del pedido. \\
\bottomrule
\end{tabularx}

\clearpage
\section{Referencias}\label{sec:referencias}
\begin{enumerate}
\item U.S. Department of Housing and Urban Development (2014). \emph{User Manual Template}, PPM 2.0. Se tomó como referencia para describir la primera sesión, las versiones, el flujo de datos y los modos de operación. \href{https://www.hud.gov/sites/documents/user_man.pdf}{Consultar la plantilla de HUD}.
\item Centers for Medicare \& Medicaid Services. \emph{Electronic Retroactive Processing Transmission User Manual}. La sección 5 orientó el diagnóstico por síntoma y los datos del reporte de incidentes. \href{https://www.cms.gov/research-statistics-data-and-systems/cms-information-technology/mapdhelpdesk/downloads/electronic-retroactive-processing-transmission-erpt-user-manual-v19.pdf}{Consultar el manual de CMS}.
\item ISO/IEC/IEEE 26514:2022. \emph{Design and development of information for users}. El resumen público orientó la organización de tareas y los resultados que cada usuario puede verificar. \href{https://www.iso.org/standard/77451.html}{Consultar la ficha de ISO}.
\item Swagger. \emph{Bearer Authentication} y demostración \emph{Petstore}. Referencias para el encabezado Bearer y el uso de \emph{Try it out}. \href{https://swagger.io/docs/specification/v3_0/authentication/bearer-authentication/}{Consultar la documentación}. \href{https://petstore.swagger.io/?displayOperationId=1}{Abrir Petstore}.
\item textbee. \emph{Getting started}. Requisitos del dispositivo Android para enviar SMS. \href{https://textbee.dev/docs/getting-started}{Consultar la documentación}.
\item Food Court. \code{docs/requirements/project-schedule.md}, \code{docs/testing/seed-data.md}, \code{scripts/seed.sh}, archivos Compose y código fuente de los cuatro servicios. Las rutas y respuestas del manual se contrastaron con los controladores, las pruebas y los contratos \code{/v3/api-docs} del entorno local.
\end{enumerate}

\clearpage
\appendix
\section{Operaciones por historia de usuario}\label{sec:rutas}
\footnotesize
\begin{minipage}{\linewidth}
\captionof{table}{Correspondencia entre historias, operaciones y roles}
\begin{tabularx}{\linewidth}{@{}P{8mm}P{12mm}P{51mm}Y P{11mm}@{}}
\toprule
HU & Puerto & Método y ruta & Rol & Éxito \\
\midrule
1 & 8081 & \route{POST}{/owners} & Administrador & 201 \\
2 & 8082 & \route{POST}{/restaurants} & Administrador & 201 \\
3 & 8082 & \route{POST}{/dishes} & Propietario & 201 \\
4 & 8082 & \route{PUT}{/dishes/{id}} & Propietario & 200 \\
5 & 8081 & \route{POST}{/login} & Público & 200 \\
6 & 8081 & \route{POST}{/employees} & Propietario & 201 \\
7 & 8082 & \route{PUT}{/dishes/{id}/active} & Propietario & 200 \\
8 & 8081 & \route{POST}{/customers} & Público & 201 \\
9 & 8082 & \route{GET}{/restaurants} & Cliente & 200 \\
10 & 8082 & \route{GET}{/restaurants/{restaurantId}/dishes} & Cliente & 200 \\
11 & 8082 & \route{POST}{/orders} & Cliente & 201 \\
12 & 8082 & \route{GET}{/orders?status=PENDING} & Empleado & 200 \\
13 & 8082 & \route{POST}{/orders/{orderId}/assignment} & Empleado & 200 \\
14 & 8082 & \route{POST}{/orders/{orderId}/ready} & Empleado asignado & 200 \\
15 & 8082 & \route{POST}{/orders/{orderId}/delivery} & Empleado asignado & 200 \\
16 & 8082 & \route{POST}{/orders/{orderId}/cancellation} & Cliente del pedido & 200 \\
17 & 8083 & \route{GET}{/orders/{orderId}/trace} & Cliente del pedido & 200 \\
18 & 8083 & \route{GET}{/restaurants/{restaurantId}/efficiency} & Propietario del restaurante & 200 \\
\bottomrule
\end{tabularx}
\end{minipage}
\normalsize

\subsection{Rutas de consulta e integración}
\begin{itemize}
\item \call{GET}{/users/{id}}{8081} devuelve datos mínimos de la cuenta. Cualquier usuario puede consultar la propia y el administrador puede consultar otras.
\item \call{GET}{/restaurants/{id}}{8082} devuelve la relación de propiedad. La consulta está disponible para el administrador o el dueño correspondiente.
\item \call{GET}{/users/{id}/phone}{8081} permite obtener el teléfono de una cuenta cliente para el flujo de aviso. Puede consultarlo un administrador, un empleado o el titular de la cuenta.
\item \call{POST}{/order-status-changes}{8083} registra cambios desde Plazoleta, con el token del cliente o del empleado según la transición. Si crea un registro, devuelve \code{201}. Si ya existe un registro idéntico, devuelve \code{200}.
\item \call{POST}{/notifications/order-ready}{8084} recibe \code{orderId}, \code{restaurantName}, \code{phone} y \code{pin} con un token de empleado. Devuelve \code{202} al aceptar la solicitud. Para operar pedidos, use \code{/orders/{orderId}/ready} desde Plazoleta.
\end{itemize}

\clearpage
\section{Recorrido desde la terminal}\label{sec:terminal}
Use Bash, \code{curl} y \code{jq}. Ejecute el bloque en una sola terminal tras cargar la semilla. El ejemplo usa el pedido 1, el restaurante 1 y las contraseñas de demostración. En la prueba del PIN incorrecto, la opción \code{-i} de \code{curl} muestra el código de estado HTTP.

\begin{lstlisting}
set -eo pipefail
USERS=http://localhost:8081
PLAZA=http://localhost:8082
TRACE=http://localhost:8083
ORDER_ID=1
RESTAURANT_ID=1

login() {
  jq -cn --arg email "$1" --arg pass Password123 \
    '{email:$email,password:$pass}' | \
    curl -fsS -H 'Content-Type: application/json' \
      --data-binary @- "$USERS/login" | jq -er .accessToken
}
PEDRO=$(login pedro.martinez@example.com)
RONALD=$(login ronald.martinez@example.com)
CARLOS=$(login carlos.ramirez@example.com)

curl -fsS -H "Authorization: Bearer $CARLOS" \
  "$TRACE/restaurants/$RESTAURANT_ID/efficiency" | jq .
curl -fsS -H "Authorization: Bearer $PEDRO" \
  "$PLAZA/orders?status=IN_PREPARATION" | jq .
curl -fsS -X POST -H "Authorization: Bearer $PEDRO" \
  "$PLAZA/orders/$ORDER_ID/ready" | jq .
curl -sS -i -H "Authorization: Bearer $PEDRO" \
  -H 'Content-Type: application/json' --data '{"pin":"000000"}' \
  "$PLAZA/orders/$ORDER_ID/delivery"
curl -fsS -H "Authorization: Bearer $PEDRO" \
  -H 'Content-Type: application/json' --data '{"pin":"482913"}' \
  "$PLAZA/orders/$ORDER_ID/delivery" | jq .
curl -fsS -H "Authorization: Bearer $RONALD" \
  "$TRACE/orders/$ORDER_ID/trace" | jq -r '.changes[].newStatus'
curl -fsS -H "Authorization: Bearer $CARLOS" \
  "$TRACE/restaurants/$RESTAURANT_ID/efficiency" | jq .
unset PEDRO RONALD CARLOS
\end{lstlisting}

Compare con los resultados de la sección \ref{sec:demo}. Los tokens quedan en variables de la terminal hasta ejecutar \code{unset} o cerrar la sesión de Bash.
\end{document}
